\section{Induction and Recursion}

In what follows, let $A$ be a class and $R$ be a class relation on $A$. We want to prove things about the members of $A$ by induction along $R$ and define things on them by recursion along $R$, and this section says what $R$ has to look like for that to work, and then does it.

\subsection{Set-Like Relations}

The property we want is that looking below any one point of $A$ only ever sees a set's worth of $A$.

\begin{boxdefinition}[Set-Like]
    We say $(A, R)$ is \textbf{set-like} iff
    \begin{align*}
        \forall y \in A, \, \setst{x \in A}{x \, R \, y}
    \end{align*}
    is a set.
\end{boxdefinition}

The example we care about most has this property for the most basic of reasons.

\begin{boxexample}[Membership]
    $(A, \in)$ is set-like for every class $A$.
\end{boxexample}

In fact, the fact in the above example is \textbf{equivalent} to the axiom scheme of comprehension (\Cref{ZFC:Compr})! Indeed, if $A = \setst{x}{\varphi\of{x, \bar{z}}}$ and $y \in A$, then
\begin{align*}
    \setst{x \in A}{x \in y} = \setst{x \in y}{\varphi\of{x, \bar{z}}}
\end{align*}
So ``membership is set-like'' is a fancy way of saying ``Comprehension''.

But not all relations are set-like.

\begin{boxnexample}
    Let $A = \OR$. Define $R$ by the following picture:
    \begin{center}
        \begin{tikzpicture}[every node/.style = {font = \small}, dot/.style = {circle, fill, inner sep = 1.5pt}]
            \node[dot, label = right:{$1$}] (1) at (0, 0) {};
            \node[dot, label = right:{$2$}] (2) at (0, 0.8) {};
            \node (d1) at (0, 1.7) {$\vdots$};
            \node[dot, label = right:{$\omega$}] (w) at (0, 2.6) {};
            \node[dot, label = right:{$\omega + 1$}] (w1) at (0, 3.4) {};
            \node (d2) at (0, 4.3) {$\vdots$};
            \node[dot, label = right:{$0$}] (0) at (0, 5.4) {};
            \draw (1) -- (2);
            \draw (2) -- (d1);
            \draw (d1) -- (w);
            \draw (w) -- (w1);
            \draw (w1) -- (d2);
            \draw (d2) -- (0);
        \end{tikzpicture}
    \end{center}
    That is, it looks a bit like $\parenth{\OR \setminus \set{0}, <}$.
\end{boxnexample}

\subsection{Proofs by Induction}

We now have another \textit{scheme} of facts. In what follows, start with $A$ a class and $R \subseteq A \times A$ a class binary relation. Assume $(A, R)$ is well-founded and set-like.

The first fact is that well-foundedness, which speaks of subsets of $A$, extends to subclasses.

\begin{boxproposition}[Minimal Members of Subclasses]
    If $S$ is a nonempty subclass of $A$, then $S$ has an $R$-minimal member.
\end{boxproposition}
\begin{proof}
    We may assume that $R \ssq A \times A$. Since $S \ne \emptyset$, we have some $z \in S$. Let $T$ be the $R$-closure of $\set{z}$. That means that we can define, by recursion on $n < \omega$, the sets
    \begin{align*}
        T_0 = \set{z}
        \qquad \qquad
        T_{n + 1} = \bigcup_{y \in T_n} \setst{x}{x \, R \, y}
    \end{align*}
    The fact that $(A, R)$ is set-like tells us that the union above is a union of sets, and therefore, a set.

    Next, put $T = \bigcup_{n < \omega} T_n$. This is the smallest $R$-closed set that has $z$ as a member.

    Now, $S \cap T$ is a nonempty subset of $A$. So it has an $R$-minimal member, say $x$. Here, using well-foundedness of $(A, R)$, it is easy to show that $x$ is $R$-minimal among members of $S$. Otherwise, we'd have $y \, R \, x$, with $y \in S$. Since $x \in T$ and $T$ is $R$-closed, we have $y \in T$. This contradicts $R$-minimality of $x$ in $S \cap T$.
\end{proof}

From this we get induction along $R$.

\begin{boxtheorem}[Theorem Scheme about Proofs by Induction]
    Let $\varphi(u, w)$ be a formula in the language of set theory, and assume $z$ is a set.

    Assume that for all $y$, if $y$ satisfies the formula that defines $A$ (that is, $\forall y \in A$),
    \begin{align*}
        \brac{\parenth{\forall x \, R \, y, \, \varphi\of{x, z}} \to \varphi\of{y, z}}
    \end{align*}
    We can then conclude that $\forall x \in A$, $\varphi\of{x, z}$.
\end{boxtheorem}
\begin{proof}
    Deny for contradiction. Let $S = \setst{x \in A}{\neg \varphi\of{x, z}}$. By the previous proposition, as $S$ is a nonempty subclass of $A$, $S$ has an $R$-minimal member, say $y$. Then, $\forall x \, R \, y$, $\varphi\of{x, z}$.

    By assumption, $\varphi\of{y, z}$, ie, $y \notin S$. This is a contradiction.
\end{proof}

\subsection{Definitions by Recursion}

Notice that in mathematics, we define things by recursion and prove things by induction. But when we explain how everything works, we usually go backwards!

\begin{boxtheorem}[Theorem Scheme about Definitions by Recursion]
    Let $F : A \times V \to V$ be a class function. Then, there is a class function $G : A \to V$ such that $\forall y \in A$,
    % [CLAUDE]
    \begin{align*}
        G(y) = F\of{y, G \restriction \setst{x}{x \, R \, y}}
    \end{align*}
    Moreover, $G$ is unique in the sense that if $H : A \to V$ is any class function such that
    \begin{align*}
        \forall y \in A,\, H(y) = F\of{y, H \restriction \setst{x}{x \, R \, y}}
    \end{align*}
    then $\forall x \in A$, $H(x) = G(x)$.
\end{boxtheorem}
\begin{boxwarning}
    Important subtlety: the statement ``there exists a class function'' is not a valid statement in the language of set theory, as we cannot quantify over \textit{classes}. What \textit{is} valid, though, is saying something like ``such and such a $G$ satisfies such and such a property'', where the definition shall come from the proof below.
\end{boxwarning}

The proof builds $G$ out of set-sized pieces, which deserve a name.

\begin{boxdefinition}[Approximation]
    % [CORRECTED] was: $g(y) = F\of{y, G \restriction \setst{x}{x \, R \, y}}$ --- $G$ is what the proof is about to build out of the approximations, so the defining equation of an approximation has to restrict $g$ itself, as every later step of the proof uses it
    We call $g$ an \textbf{approximation} iff $g$ is a set function whose domain is an $R$-closed subset of $A$ (ie, if $y \in \dom{g}$ and $x \, R \, y$, then $x \in \dom{g}$) and $\forall y \in \dom{g}$, $g(y) = F\of{y, g \restriction \setst{x}{x \, R \, y}}$.
\end{boxdefinition}

\begin{proof}
    We'll prove uniqueness and then existence.

    \begin{description}
        \item[\underline{Uniqueness.}]
            Suppose $g$ and $h$ are approximations. Then, $\forall y \in \dom{g} \cap \dom{h}$, $g(y) = h(y)$. You can prove this using $R$-induction.

        \item[\underline{Existence.}]
            We'll show that for all $z \in A$, there is some approximation $g_z$ with $\dom{g_z}$ being the $R$-closure of $\set{z}$. These $g_z$, if we can show them to exist, are unique (by above).

            Use $R$-induction. The induction hypothesis says that we have $g_y$ as required for each $y \, R \, z$. Let $h = \bigcup_{y \, R \, z} g_y$. We can see that $\dom{h}$ is the $R$-closure of $\setst{y}{y \, R \, z}$.

            Put $g_z = h \cup \set{\parenth{z, F\of{z, h \restriction \setst{x}{x \, R \, z}}}}$.

            This works! Then, put $G = \bigcup_{z \in A} g_z$.

            Indeed, $h$ is a set, by Replacement, since $y \mapsto g_y$ is a class function on the set $\setst{y}{y \, R \, z}$; it is a function, because any two of the $g_y$ agree on their common domain by uniqueness; and it is an approximation, because a union of $R$-closed sets is $R$-closed and each $g_y$ satisfies the defining equation on its own domain. Moreover, $z \notin \dom{h}$, as otherwise $z$ would lie on an $R$-cycle, which well-foundedness forbids. So $g_z$ is a function, its domain $\dom{h} \cup \set{z}$ is the $R$-closure of $\set{z}$, and it satisfies the defining equation at $z$ by construction and below $z$ because $h$ does.

            Finally, $G$ is a class function because the $g_z$ agree wherever their domains overlap, again by uniqueness, and every $y \in A$ lies in $\dom{g_y}$, so $G(y) = g_y(y) = F\of{y, g_y \restriction \setst{x}{x \, R \, y}} = F\of{y, G \restriction \setst{x}{x \, R \, y}}$, since $G$ agrees with $g_y$ on $\dom{g_y} \supseteq \setst{x}{x \, R \, y}$. Uniqueness of $G$ in the sense of the theorem is $R$-induction once more: if $H$ satisfies the same equation, then $H(x) = G(x)$ for all $x \in A$.
    \end{description}
\end{proof}

The proof also tells us how to compute $G$, which is what makes ``$G(x) = s$'' a statement in the language of set theory after all.

\begin{boxlemma}
    Given any $x \in A$ and set $s$, TFAE:
    \begin{enumerate}
        \item $G(x) = s$
        \item There exists an approximation $g$ with $x \in \dom{g}$ and $g(x) = s$.
        \item For all approximations $g$, if $x \in \dom{g}$ then $g(x) = s$.
    \end{enumerate}
\end{boxlemma}
